\chapter{Start here}\label{ch:start}

Mix and Learn is one program that does three jobs. It \emph{teaches}: lessons, a sound lab and an explanation for every control. It \emph{measures}: load your own tracks and it finds their tempo, key, loudness, brightness and structure. And it \emph{mixes}: it plans a set, explains every decision, renders the mix as an MP3 or WAV file, and saves a small set file from which the same mix can be made again.

It is a web app: plain HTML, CSS and JavaScript with no build step. On a computer it runs in any modern browser. On Android the very same files are wrapped into an app with Capacitor (Chapter~\ref{ch:android}). There is one code base, so the desktop and phone versions never drift apart.

\section{Run it on your computer}
You need Python (any recent version) to serve the files; nothing else.
\begin{lstlisting}[language=bash]
git clone https://github.com/AhmadAmmar/djmixer.git
cd djmixer
python dev/serve.py
\end{lstlisting}
Your browser opens \texttt{http://localhost:8765}. The server watches every file in \texttt{web/} and the git branch: whenever you save a file, commit, pull or rebuild, every open page reloads itself, and the server tells the browser never to use a cached copy. You always see the latest version without pressing anything.

\section{The first visit}
The first time the app opens, a short tour points at each part in turn (Figure~\ref{fig:tour}); the \emph{Tour} button repeats it. After that, anything with a dotted underline or a dashed border has an explanation: hover over it with a mouse, or press and hold it on a touch screen. Each explanation has four parts: what it is, what it means musically, the science behind it, and the file and function in the code.

\begin{figure}[H]\centering
\includegraphics[width=\linewidth]{tour.png}
\caption{The first-run tour, here on its second step, the Library. The highlighted part has an orange outline; the box at the bottom right explains it.}\label{fig:tour}
\end{figure}

\begin{figure}[H]\centering
\includegraphics[width=\linewidth]{help.png}
\caption{Hovering over \emph{Sync B to A} opens its explanation: what it does, the musical reason, the formula, and where the code is.}\label{fig:help}
\end{figure}

\section{The parts of the screen}
\begin{description}
\item[Lessons] (left) Seventeen short lessons, each with words, something to try, the formula and the code. They run from \emph{sound is air moving} to \emph{your first real mix}.
\item[Decks A and B] Two players, as on a DJ console. Each knob shows its value and has its own reset button (or double-click the slider):
\begin{itemize}
\item \textbf{Volume}, default $\times1.00$ (0\,dB): how loud this deck plays before the crossfader.
\item \textbf{Speed}, default $\times1.000$: shown also as a percentage and in semitones, because faster also means higher.
\item \textbf{Bass} (below 250\,Hz), \textbf{Mids} (around 1\,kHz) and \textbf{Treble} (above 4\,kHz), default 0\,dB each, from $-30$ to $+6$\,dB: the three equaliser bands. Bass is kick and bassline, mids are voices and melody, treble is hi-hats and the shimmer of strings.
\item \textbf{Play}: the whole track, vocals only or music only (approximate, see Chapter~\ref{ch:mixing}).
\end{itemize}
Between the decks: the crossfader, \emph{Sync B to A}, \emph{Swap A $\leftrightarrows$ B} (exchanges tracks, positions and settings), \emph{A $\to$ B} and \emph{B $\to$ A} (put the same track on the other deck), and \emph{Reset both decks}. Load a deck with its \emph{Load} button, by dropping an audio file from a folder onto it, or by dragging a track's card from the Library.

\section{One library, remembered}
Every track, however it arrives (the Library button, a drop on a deck, or a drop on the Library), becomes one record in the Library, measured once. The decks, the Library and the Set builder all point at those same records: a track dropped on deck A appears in the Library marked \emph{on deck A}, and \emph{Use the deck tracks} in the Set builder turns the two deck tracks into a set. The records, with the original file bytes and the measurements, are kept in the browser's own database (IndexedDB) on this device, so they survive reloads and app updates, and the decks reload the tracks they last held. \emph{Forget} removes a track. The memory belongs to the address: clearing the browser's data for it, or opening the app at another address or port, starts empty.

\section{Both decks together}
Below the decks, a strip shows both decks' beats on one time line, with now as the centre line: tall marks are beat one of each bar. When A's and B's marks sit on top of each other the beats are in step; the text gives the gap in milliseconds, which deck is ahead, how fast they drift apart, and each deck's volume. \emph{Play both} starts the decks on the same tick of the audio clock.

\emph{Advise the mix} (it also runs whenever a deck changes) explains how to mix A into B: where to leave A and why, what speed, start point and EQ B needs and why, and the style for the two together, with its rule. \emph{Apply to decks} sets those values; \emph{Play the transition} performs it, with every move timed on the audio clock. Under it, the best next tracks from the Library are listed with their reasons, ready to load on deck A, because DJs alternate decks. \emph{Hide lessons} (top left) gives the decks the whole width.
\item[Library] Add audio files; each is measured and drawn. A table shows how well every pair would mix.
\item[Set builder] Choose tracks, a transition style and an order, then render.
\item[Sound lab] Hear the ideas: a note, an octave, a fifth, chords, timbre, beats and bars, half time.
\item[Sargam, Spectrum, Tone tool] Live views of pitch against the swaras, of all frequencies at once, and of a single tone.
\end{description}

The \emph{Theme} button switches between light, dark and auto, which follows your computer or phone setting; the choice is remembered. Every colour, including those drawn on the waveforms and graphs, comes from the theme, so nothing disappears against the background.

Nothing you load leaves your device: files are read by the browser itself, and no audio is copied into project or set files. That matters both for privacy and because commercial music may not be shared.
